% Options for packages loaded elsewhere
% Options for packages loaded elsewhere
\PassOptionsToPackage{unicode}{hyperref}
\PassOptionsToPackage{hyphens}{url}
\PassOptionsToPackage{dvipsnames,svgnames,x11names}{xcolor}
%
\documentclass[
  letterpaper,
  DIV=11,
  numbers=noendperiod]{scrartcl}
\usepackage{xcolor}
\usepackage{amsmath,amssymb}
\setcounter{secnumdepth}{-\maxdimen} % remove section numbering
\usepackage{iftex}
\ifPDFTeX
  \usepackage[T1]{fontenc}
  \usepackage[utf8]{inputenc}
  \usepackage{textcomp} % provide euro and other symbols
\else % if luatex or xetex
  \usepackage{unicode-math} % this also loads fontspec
  \defaultfontfeatures{Scale=MatchLowercase}
  \defaultfontfeatures[\rmfamily]{Ligatures=TeX,Scale=1}
\fi
\usepackage{lmodern}
\ifPDFTeX\else
  % xetex/luatex font selection
\fi
% Use upquote if available, for straight quotes in verbatim environments
\IfFileExists{upquote.sty}{\usepackage{upquote}}{}
\IfFileExists{microtype.sty}{% use microtype if available
  \usepackage[]{microtype}
  \UseMicrotypeSet[protrusion]{basicmath} % disable protrusion for tt fonts
}{}
\makeatletter
\@ifundefined{KOMAClassName}{% if non-KOMA class
  \IfFileExists{parskip.sty}{%
    \usepackage{parskip}
  }{% else
    \setlength{\parindent}{0pt}
    \setlength{\parskip}{6pt plus 2pt minus 1pt}}
}{% if KOMA class
  \KOMAoptions{parskip=half}}
\makeatother
% Make \paragraph and \subparagraph free-standing
\makeatletter
\ifx\paragraph\undefined\else
  \let\oldparagraph\paragraph
  \renewcommand{\paragraph}{
    \@ifstar
      \xxxParagraphStar
      \xxxParagraphNoStar
  }
  \newcommand{\xxxParagraphStar}[1]{\oldparagraph*{#1}\mbox{}}
  \newcommand{\xxxParagraphNoStar}[1]{\oldparagraph{#1}\mbox{}}
\fi
\ifx\subparagraph\undefined\else
  \let\oldsubparagraph\subparagraph
  \renewcommand{\subparagraph}{
    \@ifstar
      \xxxSubParagraphStar
      \xxxSubParagraphNoStar
  }
  \newcommand{\xxxSubParagraphStar}[1]{\oldsubparagraph*{#1}\mbox{}}
  \newcommand{\xxxSubParagraphNoStar}[1]{\oldsubparagraph{#1}\mbox{}}
\fi
\makeatother


\usepackage{longtable,booktabs,array}
\newcounter{none} % for unnumbered tables
\usepackage{calc} % for calculating minipage widths
% Correct order of tables after \paragraph or \subparagraph
\usepackage{etoolbox}
\makeatletter
\patchcmd\longtable{\par}{\if@noskipsec\mbox{}\fi\par}{}{}
\makeatother
% Allow footnotes in longtable head/foot
\IfFileExists{footnotehyper.sty}{\usepackage{footnotehyper}}{\usepackage{footnote}}
\makesavenoteenv{longtable}
\usepackage{graphicx}
\makeatletter
\newsavebox\pandoc@box
\newcommand*\pandocbounded[1]{% scales image to fit in text height/width
  \sbox\pandoc@box{#1}%
  \Gscale@div\@tempa{\textheight}{\dimexpr\ht\pandoc@box+\dp\pandoc@box\relax}%
  \Gscale@div\@tempb{\linewidth}{\wd\pandoc@box}%
  \ifdim\@tempb\p@<\@tempa\p@\let\@tempa\@tempb\fi% select the smaller of both
  \ifdim\@tempa\p@<\p@\scalebox{\@tempa}{\usebox\pandoc@box}%
  \else\usebox{\pandoc@box}%
  \fi%
}
% Set default figure placement to htbp
\def\fps@figure{htbp}
\makeatother





\setlength{\emergencystretch}{3em} % prevent overfull lines

\providecommand{\tightlist}{%
  \setlength{\itemsep}{0pt}\setlength{\parskip}{0pt}}



 


\KOMAoption{captions}{tableheading}
\makeatletter
\@ifpackageloaded{caption}{}{\usepackage{caption}}
\AtBeginDocument{%
\ifdefined\contentsname
  \renewcommand*\contentsname{Table of contents}
\else
  \newcommand\contentsname{Table of contents}
\fi
\ifdefined\listfigurename
  \renewcommand*\listfigurename{List of Figures}
\else
  \newcommand\listfigurename{List of Figures}
\fi
\ifdefined\listtablename
  \renewcommand*\listtablename{List of Tables}
\else
  \newcommand\listtablename{List of Tables}
\fi
\ifdefined\figurename
  \renewcommand*\figurename{Figure}
\else
  \newcommand\figurename{Figure}
\fi
\ifdefined\tablename
  \renewcommand*\tablename{Table}
\else
  \newcommand\tablename{Table}
\fi
}
\@ifpackageloaded{float}{}{\usepackage{float}}
\floatstyle{ruled}
\@ifundefined{c@chapter}{\newfloat{codelisting}{h}{lop}}{\newfloat{codelisting}{h}{lop}[chapter]}
\floatname{codelisting}{Listing}
\newcommand*\listoflistings{\listof{codelisting}{List of Listings}}
\makeatother
\makeatletter
\makeatother
\makeatletter
\@ifpackageloaded{caption}{}{\usepackage{caption}}
\@ifpackageloaded{subcaption}{}{\usepackage{subcaption}}
\makeatother
\usepackage{bookmark}
\IfFileExists{xurl.sty}{\usepackage{xurl}}{} % add URL line breaks if available
\urlstyle{same}
\hypersetup{
  pdftitle={What Happens When Your Favorite Coffee Store Temporarily Closes?},
  colorlinks=true,
  linkcolor={blue},
  filecolor={Maroon},
  citecolor={Blue},
  urlcolor={Blue},
  pdfcreator={LaTeX via pandoc}}


\title{What Happens When Your Favorite Coffee Store Temporarily Closes?}
\usepackage{etoolbox}
\makeatletter
\providecommand{\subtitle}[1]{% add subtitle to \maketitle
  \apptocmd{\@title}{\par {\large #1 \par}}{}{}
}
\makeatother
\subtitle{Evidence from a Large Coffee Chain}
\author{}
\date{2026-04-23}
\begin{document}
\maketitle


\section{The Question}\label{the-question}

When a store closes temporarily, do customers who \emph{wanted} to buy
but couldn't --- those with a purchase blocked --- break their coffee
habit? Or do they come back as if nothing happened?

The standard intuition is that \textbf{habits are fragile}: if a planned
purchase is missed, the routine weakens and you buy less afterward. We
test this idea using data from a large coffee chain, where 22 stores
closed unexpectedly for 10--29 days. We focus specifically on customers
who had a purchase due during the closure --- those for whom the closure
actually blocked an intended purchase --- and track whether their
behavior changed after reopening.

We study two behaviors:

\begin{itemize}
\tightlist
\item
  \textbf{How much they buy} --- did closure reduce purchase frequency?
\item
  \textbf{What they buy} --- did closure change whether they try new
  products?
\end{itemize}

\begin{center}\rule{0.5\linewidth}{0.5pt}\end{center}

\section{Setting and Data}\label{setting-and-data}

The coffee chain is a Chinese quick-service brand with millions of
customers ordering primarily through a mobile app. Our data cover June
2020 -- December 2021 and include 10.6 million transactions from 779,000
customers.

We identify \textbf{22 store closures} (10--29 days long) and assign
customers to two groups based on which store they visited most before
the closure:

\begin{itemize}
\tightlist
\item
  \textbf{Treatment group} --- customers whose preferred store closed
  (\textasciitilde8,000 customers)
\item
  \textbf{Control group} --- customers whose preferred store stayed open
  (\textasciitilde37,000 customers)
\end{itemize}

\begin{center}\rule{0.5\linewidth}{0.5pt}\end{center}

\section{The Key Distinction: Blocked Buyers vs.~Non-Blocked
Buyers}\label{the-key-distinction-blocked-buyers-vs.-non-blocked-buyers}

Not every treatment customer would have bought during the closure even
if the store had stayed open. Only frequent buyers --- those with a
\textbf{purchase due} during the closure window --- were truly blocked
by the closure. We call these \textbf{blocked buyers}.

We use a machine learning model (XGBoost) trained on pre-closure
purchasing patterns to classify each treatment and control customer as a
``blocked buyer'' (would have bought during the closure) or
``non-blocked buyer'' (would not have bought regardless). The model
achieves \textasciitilde73\% accuracy.

This distinction is crucial: \textbf{habit disruption, if real, should
show up most strongly among blocked buyers} --- they are the ones whose
routine was actually interrupted.

\begin{center}\rule{0.5\linewidth}{0.5pt}\end{center}

\section{Research Design}\label{research-design}

We use a \textbf{Triple Difference-in-Differences (DDD)} design that
simultaneously compares:

\begin{itemize}
\tightlist
\item
  Treatment vs.~control customers
\item
  Before vs.~after the closure
\item
  Blocked buyers vs.~non-blocked buyers
\end{itemize}

This lets us separately estimate two effects:

{\def\LTcaptype{none} % do not increment counter
\begin{longtable}[]{@{}
  >{\raggedright\arraybackslash}p{(\linewidth - 2\tabcolsep) * \real{0.5000}}
  >{\raggedright\arraybackslash}p{(\linewidth - 2\tabcolsep) * \real{0.5000}}@{}}
\toprule\noalign{}
\begin{minipage}[b]{\linewidth}\raggedright
Effect
\end{minipage} & \begin{minipage}[b]{\linewidth}\raggedright
What it measures
\end{minipage} \\
\midrule\noalign{}
\endhead
\bottomrule\noalign{}
\endlastfoot
\textbf{Baseline effect} (\(\delta^B\)) & Post-closure change for
\emph{non-blocked} treatment customers relative to control --- captures
general demand shifts from the closure (e.g., discovering alternatives,
post-reopening promotions) \\
\textbf{Blocked-buyer effect} (\(\delta^D\)) & \emph{Additional}
post-closure change for blocked treatment customers, over and above the
baseline --- this is the causal effect of having been \textbf{blocked
from a planned purchase} \\
\end{longtable}
}

The regression equation:

\begin{align}
Y_{iet} &= \delta^B \cdot \text{post} \times \text{Treated}
         + \beta \cdot \text{post} \times \text{Blocked} \\
        &\quad + \delta^D \cdot \text{post} \times \text{Treated} \times \text{Blocked}
         + \phi_{ie} + \omega_t + \gamma_m + \varepsilon_{iet}
\end{align}

where \(\phi_{ie}\) = customer--closure fixed effects, \(\omega_t\) =
relative-period fixed effects, \(\gamma_m\) = calendar-month fixed
effects. Standard errors are clustered at the member level.

The four-group prediction (post-closure change in outcome):

{\def\LTcaptype{none} % do not increment counter
\begin{longtable}[]{@{}lll@{}}
\toprule\noalign{}
& Non-blocked & Blocked \\
\midrule\noalign{}
\endhead
\bottomrule\noalign{}
\endlastfoot
\textbf{Control} & 0 & \(\beta\) \\
\textbf{Treatment} & \(\delta^B\) & \(\delta^B + \beta + \delta^D\) \\
\end{longtable}
}

\begin{center}\rule{0.5\linewidth}{0.5pt}\end{center}

\section{Results}\label{results}

\subsection{Purchase Frequency}\label{purchase-frequency}

Did closures reduce how often customers buy at the chain?

{\def\LTcaptype{none} % do not increment counter
\begin{longtable}[]{@{}lrrl@{}}
\toprule\noalign{}
Coefficient & Estimate & SE & \\
\midrule\noalign{}
\endhead
\bottomrule\noalign{}
\endlastfoot
post × treated (\(\delta^B\)) & +0.00081 & 0.00064 & \\
post × blocked (\(\beta\)) & −0.03353 & 0.00107 & *** \\
post × treated × blocked (\(\delta^D\)) & +0.00517 & 0.00241 & ** \\
\end{longtable}
}

\emph{358,864 observations; } p\textless0.1, ** p\textless0.05, ***
p\textless0.01.*

\textbf{What this means:}

\begin{itemize}
\tightlist
\item
  \textbf{No general habit disruption} (\(\delta^B \approx 0\)):
  Treatment customers who were \emph{not} blocked buy just as much as
  before --- the closure itself caused no demand loss.
\item
  \textbf{Blocked buyers regress to the mean} (\(\beta = -0.034\)):
  Blocked buyers (in both treatment and control) are high-frequency
  purchasers before the closure and naturally buy less afterward. This
  is \textbf{not} caused by the closure --- it happens equally in the
  control group and reflects a normal statistical mean-reversion for
  heavy buyers.
\item
  \textbf{No habit disruption for blocked buyers either}
  (\(\delta^D = +0.005\)): The blocked-buyer effect is small and
  \emph{positive}, meaning blocked treatment customers recovered at
  least as fast as blocked control customers. There is no evidence the
  forced missed purchase reduced future buying.
\end{itemize}

A dynamic event-study confirms no pre-trends (pre-period joint test:
\(p = 0.907\)) and shows a \textbf{short-run rebound} in purchase
frequency immediately after reopening (significant in weeks 1 and 3),
which then fades --- consistent with pent-up demand being satisfied, not
with any persistent disruption.

\subsection{Novelty-Seeking}\label{novelty-seeking}

Did closures change what kinds of products customers try?

{\def\LTcaptype{none} % do not increment counter
\begin{longtable}[]{@{}lrrl@{}}
\toprule\noalign{}
Coefficient & Estimate & SE & \\
\midrule\noalign{}
\endhead
\bottomrule\noalign{}
\endlastfoot
post × treated (\(\delta^B\)) & +0.031 & 0.011 & *** \\
post × blocked (\(\beta\)) & +0.036 & 0.006 & *** \\
post × treated × blocked (\(\delta^D\)) & −0.039 & 0.014 & *** \\
\end{longtable}
}

\emph{358,864 observations; } p\textless0.1, ** p\textless0.05, ***
p\textless0.01.*

\textbf{What this means:}

\begin{itemize}
\tightlist
\item
  \textbf{Non-blocked treatment customers explore more}
  (\(\delta^B = +0.031\)): When routine is disrupted but no specific
  purchase was blocked, customers become more adventurous --- they are
  more likely to try products they have never bought before. Closure
  breaks the autopilot.
\item
  \textbf{Blocked buyers are more exploratory in general}
  (\(\beta = +0.036\)): Heavy buyers encounter the menu more frequently
  and naturally try more new items, regardless of closure.
\item
  \textbf{Blocked treatment customers explore less}
  (\(\delta^D = -0.039\)): This is the key finding. Among customers who
  had a purchase blocked, treatment customers try significantly
  \emph{fewer} new products after reopening compared to equally-frequent
  control buyers. The blocked purchase goal narrows their attention:
  they return to fulfill the specific, familiar purchase they missed ---
  not to explore.
\end{itemize}

\begin{center}\rule{0.5\linewidth}{0.5pt}\end{center}

\section{Why Does This Happen? Two
Mechanisms}\label{why-does-this-happen-two-mechanisms}

\subsection{Mechanism 1: Blocked Goals Create Focused
Rebound}\label{mechanism-1-blocked-goals-create-focused-rebound}

When a purchase is blocked, the unfulfilled intention does not simply
disappear. Research in psychology shows that \textbf{unfinished goals
stay mentally active} and generate motivation to complete them (the
Zeigarnik effect). Combined with \textbf{psychological reactance} ---
the desire to reclaim a restricted freedom --- this creates a focused
drive to return and fulfill the specific original purchase.

This explains both findings simultaneously: - Blocked buyers buy
\emph{more} immediately after reopening (goal completion drive →
purchase frequency rebound) - Blocked buyers explore \emph{less}
(narrowed focus on completing the specific familiar purchase → reduced
novelty-seeking)

Supporting evidence: 7.3\% of blocked treatment customers purchased at a
\emph{different} store of the same chain during the closure --- direct
evidence that the blocked goal was actively pursued through substitution
rather than abandoned.

\subsection{Mechanism 2: Missed Exposure to New
Products}\label{mechanism-2-missed-exposure-to-new-products}

During the closure, control customers continue buying and encounter
newly introduced menu items. Blocked treatment customers miss this
exposure window entirely. After reopening, they lack familiarity with
the updated assortment and default to previously known products ---
reducing their novelty rate relative to control buyers who experienced
the same menu evolution.

Descriptive evidence: matched control stores introduced an average of
\textbf{7.9 new products} during each closure window (range: 1.0--15.4).
Blocked treatment customers missed exposure to all of them.

A direct test: if this mechanism is operating, the negative
blocked-buyer effect on novelty-seeking (\(\delta^D\)) should be larger
at closures where control stores introduced more new products. This is a
planned robustness check.

\begin{center}\rule{0.5\linewidth}{0.5pt}\end{center}

\section{Summary}\label{summary}

{\def\LTcaptype{none} % do not increment counter
\begin{longtable}[]{@{}
  >{\raggedright\arraybackslash}p{(\linewidth - 2\tabcolsep) * \real{0.5000}}
  >{\raggedright\arraybackslash}p{(\linewidth - 2\tabcolsep) * \real{0.5000}}@{}}
\toprule\noalign{}
\begin{minipage}[b]{\linewidth}\raggedright
Finding
\end{minipage} & \begin{minipage}[b]{\linewidth}\raggedright
Result
\end{minipage} \\
\midrule\noalign{}
\endhead
\bottomrule\noalign{}
\endlastfoot
Did closures reduce purchase frequency for blocked buyers? & \textbf{No}
--- blocked buyers recover at least as fast as control buyers with
identical purchase frequency \\
Did closures increase exploration of new products overall? &
\textbf{Yes} --- non-blocked treatment customers try more new products
after reopening \\
Did blocked buyers explore more or less than equally-frequent control
buyers? & \textbf{Less} --- their blocked purchase goal narrows
attention to familiar products \\
Is the null purchase-frequency result consistent with habit disruption?
& \textbf{No} --- if anything, blocked buyers bounce back slightly
faster, not slower \\
\end{longtable}
}

\textbf{The big picture:} Among customers who had a purchase blocked by
the closure, the habit is not eroded --- they return to buying at the
same rate as equally-frequent buyers whose store never closed. What
changes is \emph{what} they buy: a blocked purchase goal sends customers
back on a focused mission to complete the specific purchase they missed,
leaving them less open to trying new products than comparable buyers who
were never interrupted. For customers without a blocked purchase goal,
the effect runs in the opposite direction --- disrupted routine opens
them up to explore the menu more broadly.




\end{document}
